% ---------------------------------------------------------------------------
% Appendix: Experimental Details -- regenerated for the cluster_full run.
% Every number below is read from configs/cluster_full/config.yaml,
% configs/cluster_full/ATTACK_SPECS.md and scripts/cluster/config.sh.
% Requires: booktabs, amsmath.
% ---------------------------------------------------------------------------
\section{Experimental Details}
\label{app:exp_details}

\paragraph{Attacked models and data.}
The attacked model is, for each dataset, the public Riemannian flow-matching checkpoint of
\citet{dufour2024around} on $\mathcal{S}^2$ (\texttt{PLONK\_YFCC}, \texttt{PLONK\_OSV\_5M}).
The image encoder $\mathcal{E}$ is frozen and shared with the diffusion and Euclidean
flow-matching checkpoints used in the transfer experiment: DINOv2 ViT-L/14 with registers
\citep{oquab2024dinov2} on YFCC4k, the StreetCLIP ViT-L/14 vision tower
\citep{haas2023streetclip} on OSV-5M, both taking a $336{\times}336$ input obtained by a
$1{:}1$ centre crop, bicubic resize and the encoder's own channel normalisation. We evaluate
$4\,000$ images per dataset: YFCC4k in full, and for OSV-5M a fixed, prefix-stable subset of
the test split (seeded shuffle, seed $42$), so that every attack, budget and ablation sees
exactly the same images.

\paragraph{Attack hyperparameters.}
Attacks are trained per image over 60 PGD steps, with batch-size 256 (where applicable) and
learning rate $10^{-3}$. These hyperparameters ensure comparable training length and memory
usage across considered baselines despite their design differences. We train with budgets
$b \in \{1/255,\, 4/255,\, 8/255\}$. This choice aims at demonstrating the high performance of
our attack even at very low perturbation budgets, considered as highly stealthy in the
literature \citep{kang2024diffattackevasionattacksdiffusionbased,
croce2021robustbenchstandardizedadversarialrobustness}. Budget $8/255$ allows direct
comparison with GeoShield. Concretely, every trainable attack shares the same outer loop: the
perturbation $\delta$ lives on the $[0,1]$ pixel scale of the model-input square and is
updated by sign-SGD, $\delta \leftarrow \mathrm{clip}(\delta - \eta\,\mathrm{sign}(\nabla_\delta
\mathcal{L}), \pm b)$ with $\eta = 10^{-3}$, from a small random start
$\delta_0 \sim \mathcal{U}(-10^{-4}, 10^{-4})$. Expectations over $(x_t, \epsilon, t)$ are
estimated with $256$ Monte-Carlo samples per step, $t \sim \mathcal{U}[0,1]$ and
$\gamma(t)$ the model's own noise schedule. The forward process is anchored at the image's
\emph{true} location $x_0$, which the defender knows by construction; only GeoFA-ACE, which is
targeted at an image rather than at a location, anchors it instead on a bank of $512$ samples
drawn once from the clean model ($100$ steps, $w=0$) and shared by all restarts.

\paragraph{Restarts and selection.}
Each attack is trained from 5 independent random initializations (except GeoShield), and we
select the one achieving the highest displacement from the true location in order to improve
attack consistency. We emphasize that this selection is not an evaluation artifact but a
\emph{deployable defender-side verification protocol}: in our setting the image's location is
the defender's own private information, so verifying candidate perturbations against the
white-box model and keeping the best is a legitimate step of the defense. Verification is also
necessary: we observe that GeoFA-TD's adversarial loss is imperfectly correlated with the
resulting FSD, so loss-based selection would be inadequate. For comparability with single-run
protocols, we report single-restart results. A restart is scored exactly like the final
evaluation, but with $128$ samples and noise seeded by $42 + \text{restart index}$, so that
the selection never reuses the evaluation noise.

\paragraph{Evaluation protocol.}
The protected image $\mathrm{clamp}(c+\delta)$ is quantised to 8 bits and fed to the model
through the very same preprocessing as a clean image, so that nothing of the perturbation
survives outside the pixel domain. For every (image, attack, budget) we draw $128$ locations
with $250$ sampling steps and guidance $w=10$ from noise shared with the clean run; FSD is the
mean haversine distance between these predictions and the ground truth, and
$\mathrm{ASR}_{>d}$ the fraction of images whose FSD exceeds $d \in \{200, 750, 2500\}$\,km.
Clean rows apply the same protocol to the unprotected image.

\paragraph{Attack-specific settings.}
Table~\ref{tab:hyperparams-specific} lists what differs between instances; everything else
(budget, optimiser, steps, restarts, evaluation) is shared. Two additional variants are
reported in the appendix only: \emph{GeoFA-Sampling} ($w{=}10$), identical to GeoFA-Sampling
except that the unrolled sampler runs at the evaluation guidance, which tests whether matching
the guidance at training time helps; and \emph{GeoFA-TD-blur}, which mixes the GeoFA-TD loss on
the protected image and on a Gaussian-blurred copy ($\sigma \sim \mathcal{U}[0,4]$, equal
weights), a blur-EOT variant meant for the robustness ablation.

\begin{table}[h]
\centering\small\setlength{\tabcolsep}{5pt}
\caption{Attack-specific settings. $\psi$ is the velocity network, $\mathcal{E}$ the image
encoder, $c$ the clean image and $x_0$ its true location. ``Gaussian'' denotes the interpolant
$x_t = \sqrt{\gamma}\,x_0 + \sqrt{1-\gamma}\,n$, $n \sim \mathcal{N}(0, I_3)$; ``geodesic'' the
forward process of Eq.~\ref{equ:forward} between a uniform point of $\mathcal{S}^2$ and the
data endpoint.}
\label{tab:hyperparams-specific}
\begin{tabular}{llll}
\toprule
Attack & Objective & Data endpoint / target & Forward process \\
\midrule
GeoFA-TD       & $\min$ cosine$\big(\psi(x_t|c+\delta), \psi(x_t|c)\big)$       & $x_0$ (true)            & Gaussian \\
GeoFA-$L^2$    & $\max \lVert \psi(x_t|c+\delta) - \psi(x_t|c) \rVert^2$        & $x_0$ (true)            & Gaussian \\
GeoFA-DPA      & $\min$ RFM loss towards $-x_0$                                 & antipode of $x_0$       & geodesic \\
GeoFA-Sampling & $\max$ haversine$\big(\hat{x}_0^{c+\delta}(x_N), x_0\big)$     & $x_0$ (true)            & $16$-step sampler, $w{=}0$ \\
GeoFA-Encoder  & $\max \lVert \mathcal{E}(c+\delta) - \mathcal{E}(c) \rVert^2$  & --                      & -- \\
GeoFA-AdvDM    & $\max$ RFM loss at $x_0$                                       & $x_0$ (true)            & geodesic \\
GeoFA-UniDef   & same as AdvDM, FDJE gradient ($h = 10^{-2}$, one probe)        & $x_0$ (true)            & geodesic \\
GeoFA-ACE      & $\min \lVert \psi(x_t|c+\delta) - \psi(x_t|c') \rVert^2 + \alpha \lVert \mathcal{E}(c+\delta) - \mathcal{E}(c') \rVert^2$ & target image $c'$, $\alpha = 1$ & geodesic, clean-model bank \\
\bottomrule
\end{tabular}
\end{table}

\paragraph{GeoShield.}
GeoShield \citep{liu2025geoshield} is run with the authors' released implementation and its
\texttt{ensemble\_3models} configuration: $100$ FGSM-style sign steps of size $1/255$ at the
budgets $\{1, 4, 8\}/255$, a CLIP ensemble (ViT-B/16, ViT-B/32, LAION ViT-G/14), global and
random-crop (scale $[0.5, 0.9]$) feature-similarity losses and a geo-text term. To make the
comparison exact, it is given lossless PNG copies of the very square the attacked pipeline
feeds to its encoder ($336$\,px), so that its own crop and resize are a no-op and its
perturbation is measured against the same clean tensor as the trainable attacks; its output is
then scored by the protocol above. We note that the released code ships a placeholder image
description and no detector regions, so our GeoShield is its CLIP-ensemble feature attack with
a constant geo-text term rather than the full VLM/detector pipeline of the original paper. It
has no restart mechanism.

\paragraph{Ablations.}
All ablations reuse the selected perturbations; nothing is retrained. \emph{Sampling steps:}
the protected image is re-evaluated with $\{8, 16, 32, 64, 250\}$ steps at $w=10$.
\emph{Guidance:} the same with $w \in \{1, 2, 5, 10\}$. \emph{Robustness:} the protected image
is JPEG-compressed at quality $\{20, 40, 60, 80, 100\}$ ($100$ = no compression) or blurred
with a Gaussian kernel of $\sigma \in \{0, 2, 4, 6, 8, 10\}$ pixels before the encoder, then
evaluated at $250$ steps; this sweep covers every attack of the run. \emph{Transfer:} the same
protected image is evaluated against the diffusion and Euclidean flow-matching checkpoints of
\citet{dufour2024around}, which share the encoder, at $250$ steps and $w=10$.
\emph{Restarts:} success rates are recomputed when keeping the best of the first $k \le 5$
restarts. \emph{Localizability:} the entropy of Eq.~\ref{equ:localizability} is estimated on
the clean image with $256$ Monte-Carlo samples of the model likelihood, and images are split
at the $33\%$ and $66\%$ quantiles.

\paragraph{Compute.}
The evaluation is sharded on a SLURM cluster as one job per (attack, window of $100$ images),
i.e. $11 \times 40 = 440$ single-GPU jobs per dataset, each covering the three budgets, the
five restarts and the ablations; the shards are then merged into the reported results. Cost is
dominated by the number of passes through $\psi$ per PGD step: GeoFA-Encoder needs none,
GeoFA-AdvDM, GeoFA-DPA and GeoFA-UniDef one forward-backward, GeoFA-TD, GeoFA-$L^2$ and
GeoFA-ACE two (the reference prediction is recomputed at every step), and GeoFA-Sampling
unrolls the $16$-step sampler.
